\section{Limitations}
\label{sec:limits}

\paragraph{No per-case ground truth.} The central limitation: no public dataset pairs cervical MRI with
per-case expert measurements, so we report distribution separation and precision, never per-case
sensitivity/specificity or absolute accuracy (e.g.\ an absolute Cobb MAE). Closing this requires an
institutional read (e.g.\ AUBMC), which we did not have.

\paragraph{Sample size and selection.} The validation used 12 healthy and 10 symptomatic necks, and the
10 unhealthy cases were \emph{selected to carry mid-cervical lesions}. Part of the Group 3 sharpness is
therefore by construction; a random draw from the 250 MMCSD cases is the necessary next test, and may
soften the separation.

\paragraph{Spectrum and cohort bias.} MMCSD is a surgical (severe-end) cervical-spondylosis cohort, so
claims are scoped to the cervical-spondylosis spectrum; generalization to mild disease or to
non-degenerative pathology is not established. The healthy controls are young (wide canals), so the
clean (0\%) behaviour of the SAC$<3$/canal$<10$ cuts on healthy may not hold for an older clinical
``normal.''

\paragraph{Findings the system does not measure.} Because measurements are geometric/signal detectors, a
pathology that does not change the measured anatomy is missed --- e.g.\ a tumour mass without geometric
distortion, foraminal stenosis (not measurable on sagittal MRI), or inflammatory signal we do not screen
for. These are scope limitations, stated explicitly; tumour (5.3) and scar (5.4) are out of scope.

\paragraph{The compression-fracture arm is under-tested.} Spondylosis does not exercise the Group 1/5.2
compression axis, and no labelled cervical compression-fracture MRI dataset surfaced in repeated hunts;
the existing 5.2 ``unhealthy'' evidence is degenerative (Duke DCM) and CT (RSNA, mostly non-compression).
This axis remains the weakest and is a documented data gap.

\paragraph{Threshold provenance and modality.} Several thresholds are borrows from radiograph/CT/lumbar
sources (Torg, some canal cuts, lumbar disc rules); they are flagged with caveats but their MRI
soft-tissue calibration is not established. The disc-degeneration grade cut-points lack published
$\kappa$/AUC for cervical and are research-grade. Spondylolisthesis has no supine-MRI presence threshold
and is reported as experimental.

\paragraph{Group 2 not validated; Group 6 not end-to-end run.} The disc group has a confirmed,
root-caused bug pending a fix with its author; the interpretation layer is built and unit-tested but its
end-to-end run is gated on the measurement validation reported here.

\paragraph{Slice thickness.} Cervical T2 is typically \SIrange{3}{4}{\milli\meter} through-plane, so the
\SI{1}{\milli\meter} isotropic measurement grid is interpolated; the system surfaces a
resolution-quality flag, but coarse acquisitions reduce confidence in fine measurements.
